\section{Introduction}\label{sec:introduction}

For a finite undirected graph $G$ on $[N]$, with nonnegative rational
capacities $c_{ij}=c_{ji}$ and $c_{ii}=0$, define
\begin{equation}\label{eq:uniform-objective}
 \Phi(G)=\min_{\varnothing\ne S\subsetneq[N]}
 \frac{\displaystyle\sum_{i\in S,\,j\notin S}c_{ij}}
 {|S|(N-|S|)}.
\end{equation}
The denominator is the number of unordered pairs separated by the cut.
Thus Equation~\eqref{eq:uniform-objective} is the sparsest-cut objective
with demand exactly one on every pair of distinct vertices.  Capacities
and thresholds are encoded in binary.  A factor-$C$ approximation must
return a nonempty proper cut whose ratio is at most $C\Phi(G)$.

\begin{theorem}\label{thm:main}
For every fixed real constant $C>1$, there is a polynomial-time reduction
from $3$CNF satisfiability to finite undirected graphs with nonnegative
rational capacities and exactly unit demands, together with a positive
rational threshold $a$, such that
\[
 \begin{array}{ll}
 \text{the formula is satisfiable}&\Longrightarrow\ \Phi(G)\le a,\\[2pt]
 \text{the formula is unsatisfiable}&\Longrightarrow\ \Phi(G)>Ca.
 \end{array}
\]
The graph size and the bit lengths of all output numbers are polynomial
in the formula size, with constants and exponents allowed to depend on
$C$.  In particular, approximating Uniform Sparsest Cut within any fixed
factor $C>1$ is NP-hard.
\end{theorem}

The source decision problem is NP-complete
\cite[Main Theorem, Problem~11]{Karp1972}.
The reduction uses no unproved complexity hypothesis.  It gives an
explicit gap for the objective in Equation~\eqref{eq:uniform-objective},
including every nonempty proper cut, however small either side may be.

\subsection{Historical context and significance}

Uniform Sparsest Cut asks for a partition that separates little capacity
relative to the number of vertex pairs it separates. This normalization
makes the problem sensitive to bottlenecks at every scale, including cuts
with a small side. Its connection with multicommodity flow gives both a
way to certify that a graph has no sparse cut and a route to approximation
algorithms. Leighton and Rao developed this flow--cut approach and its
logarithmic approximation guarantees \cite{LeightonRao1999}. Arora, Rao,
and Vazirani improved the approximation factor to $O(\sqrt{\log N})$
using semidefinite programming; their treatment of the unweighted uniform
objective appears in \cite[Section~6, Theorem~16]{ARV2009}.

Exact hardness and approximation hardness address different questions.
Matula and Shahrokhi established NP-hardness of finding sparsest cuts
\cite{MatulaShahrokhi1990}; Bonsma, Broersma, Patel, and Pyatkin proved
NP-completeness for unweighted graphs with the cardinality-product
denominator used here \cite[Theorem~15]{BBPP2012}.
To rule out a constant-factor approximation, one must create a
multiplicative separation between satisfiable and unsatisfiable instances,
while preserving the uniform denominator.

Previous approximation lower bounds depend on the complexity assumption
and on which pairs carry demand. Amb\"uhl, Mastrolilli, and Svensson proved
that a polynomial-time approximation scheme for Uniform Sparsest Cut
would yield probabilistic SAT algorithms of time $2^{n^\eps}$ for every
fixed $\eps>0$, where $n$ is the SAT input length
\cite[Theorem~1.3]{AMS2011}. Raghavendra and Steurer introduced the
Small-Set Expansion hypothesis and its connection with Unique Games
\cite{RaghavendraSteurer2010}. Raghavendra, Steurer, and Tulsiani
established arbitrarily large constant-factor gaps for Balanced Separator
from that hypothesis
\cite[full version, Theorem~3.5 and Corollary~3.6]{RST2012}.
Their balanced-cut statement measures set size by stationary volume.
The standard recursive reduction then yields hardness of every constant
factor for Uniform Sparsest Cut under polynomial-time oracle reductions,
assuming the same hypothesis.\footnote{The recursive procedure in
\cite[Section~3.3]{LeightonRao1999} calls sparsest-cut approximation on
induced subgraphs while retaining the original vertex masses. For fixed
parameters, the masses in the construction of \cite[Section~6]{RST2012}
are products of a uniform input measure and fixed rational probabilities.
Proportional replication, with sufficiently large capacities inside each
cluster, realizes these masses by polynomially many vertices with unit
all-pairs demand; compare \cite[Section~2]{ARV2009}. The resulting
constant loss is absorbed by the arbitrarily large balanced-cut gap.}

The nonuniform problem permits a selected set $\cD$ of distinct demand
pairs. In the formulation of Chawla, Krauthgamer, Kumar, Rabani, and
Sivakumar, a solution is an edge set $M$ with total capacity $c(M)$, and
its ratio is $c(M)/q_\cD(M)$, where $q_\cD(M)>0$ counts the demand pairs
disconnected after deleting $M$.
They proved every-constant-factor hardness under the Unique Games
Conjecture \cite[Corollary~1.3 and p.~99]{CKKRS2006}.
The independent Unique Games Theorem supplies the required game premise:
its unweighted bipartite constraints can be given equal weights summing
to one \cite[Theorem~1.1 and Corollary~8.3]{OpenAIUniqueGames2026}.
The resulting nonuniform hardness concerns cutset-output search under
polynomial-time Cook reductions. Only the selected pairs carry unit
demand; they need not include all distinct vertex pairs.
The direct reduction proving Theorem~\ref{thm:main} is independent of
this nonuniform consequence.

Khot and Vishnoi also developed the connection between Unique Games,
cut hardness, and negative-type metrics \cite{KhotVishnoi2015}.
Ordinary NP-hardness for nonuniform sparsest cut is known even with
treewidth-two supply graphs, with a $17/16-\eps$ approximation threshold
\cite[Theorem~1.3]{GuptaTalwarWitmer2013}. The separation between the
uniform and nonuniform
hardness questions is discussed in
\cite[Section~1.2]{ManurangsiTrevisan2018} and in the recent work of
d'Orsi, Jones, Ruotolo, Vadhan, and Zhang on low-degree Abelian Cayley
graphs \cite{DOrsiEtAl2025}.

Theorem~\ref{thm:main} establishes an unconditional approximation gap with
unit demands.  Its force comes from controlling all cut masses in a
polynomial-size reduction.  The proof starts with an explicit algebraic
test system for $3$CNF, extracts a Boolean valuation from a rare successful
event, and forces that event from any cut whose comparison cost is small
relative to its demand variance.

\subsection{Proof structure and main constructions}

The reduction has four interfaces. A satisfying assignment supplies a
low-cost cut; conversely, a cut with small cost relative to its separated
demand will determine a satisfying assignment. The main difficulty in
this converse is that the separated demand can be arbitrarily small.
Errors must therefore be controlled relative to that demand, rather than
by one fixed additive tolerance.

\paragraph{Algebraic testing.}
A proof is a bit assignment on an explicitly constructed finite set of
positions. A test samples a bounded list of positions and checks their
answers. Section~\ref{sec:pencil} constructs a linear-size list of
evaluation vectors on which a quadratic form that is nonzero somewhere
on the list vanishes only rarely. It uses the Garcia--Stichtenoth tower
\cite{GS1996,Stichtenoth2001}; the local algebra, dimension estimates,
and finite computation needed here are supplied in that section.
Section~\ref{sec:bit-tests} uses this list to construct tests whose
individual query probabilities are bounded by a fixed multiple of uniform
measure. Its auxiliary parity-query law has exactly uniform marginals
and strongly mixing pairs. It encodes Boolean tables by matrix minors and
corrects a nearly passing oracle so that quadratic line identities hold
at every base on a high-mass set of directions. Local proof testing
belongs to the PCP framework developed by Arora and Safra and by Arora,
Lund, Motwani, Sudan, and Szegedy, and subsequently recast through gap
amplification by Dinur \cite{AroraSafra1998,ALMSS1998,Dinur2007}.
The correlated test laws and their soundness are proved here.

A \emph{predicate} is a Boolean function of the entire proof assignment.
A \emph{valuation} assigns a bit to each actual predicate, so two sampled
descriptions of the same function must receive the same answer.
Calibration prescribes its mean under each specified predicate-sampling
law, and order preservation
asks that it usually respect pointwise inclusions of predicates.
Section~\ref{sec:events} proves that an approximately calibrated,
approximately order-preserving valuation would decode a satisfying
assignment.

\paragraph{Rare-event extraction.}
Section~\ref{sec:entropy} studies a Boolean response to a conjunction of
many independently sampled predicates. Its success probability may be
very small. Conditioning on success and exposing a common prefix
produces a density with nearly extremal entropy. Comparisons with the
constant-one predicate bound the density from above in mean. An explicit
clipping argument then forces it near two values, zero and a fixed
positive constant. One threshold gives a single valuation satisfying all
required calibrations and comparisons. The bounds depend on the fixed
list of predicate types, not on their support sizes or on a smallest atom.

\paragraph{From cut variance to a rare successful event.}
A \emph{score} is a bounded real function of boundedly many proof bits.
A candidate cut is a measurable $\{0,1\}$-valued function of scores.
A demand law
measures the variance of this color, and comparison measures charge for
disagreement between pairs of scores. Sections~\ref{sec:scores}
and~\ref{sec:paths} construct these measures. The demand score is a sum
of rare predicates with independent real amplitudes at several scales.
Small resampling cost forces the cut color to retain a signal from at
least one amplitude. Smoothed comparison paths preserve enough of this
signal to produce the successful conjunction response required above.
The conditional experiments ensure that this response depends on the
predicate itself. Every error is proportional to the cut variance, so
the argument applies to every positive demand mass.

\paragraph{Exact uniform demands.}
Section~\ref{sec:finite-graph} rounds each score to a grid-valued truth table. A \emph{key} consists
of that table together with its ordered essential bit names. Equal
functions have the same key. Box subdivision gives rational upper
capacities and rational lower demand masses. Replicating keys converts
these retained masses into vertex counts; heavy edges prevent a cut of
small ratio from separating copies. Additional edges from the remaining keys charge for their possible
contribution to separated demand. This produces exactly the denominator
$|S|(N-|S|)$ with polynomially many vertices and polynomial binary
encodings.

\subsection{Conventions and parameter dependence}

All probabilities refer to the explicitly described experiments.
A predicate is a Boolean function of the entire proof-bit assignment;
its sampling data may be retained when conditioning, but a response to
the predicate depends only on that function.  The same convention applies
to scores.  Thus two presentations of the same function cannot receive
different cut colors.

For a probability law $\mu$ and a Boolean function $h$,
\[
 \Var_\mu(h)=\mu(h=1)\mu(h=0)
 =\tfrac12\E_{B,B'\sim\mu}\abs{h(B)-h(B')},
\]
where the last two draws are independent.  This identity explains the
variance normalization used before the final graph is formed.

The approximation factor is fixed first.  Thereafter every numerical
constant, field size, arity, and number of repetitions is fixed independently
of the input formula.  Parameters chosen later may depend on earlier
parameters only; the necessary order is specified in the proofs.  Bounds
written $O(\cdot)$ have constants depending only on parameters already
fixed at that point.  The construction is polynomial in the formula size;
no claim is made that its exponent is uniform in $C$.
